\documentclass[11pt]{article}
\usepackage[a4paper,margin=2cm]{geometry}
\usepackage{amsmath,amssymb}
\usepackage[hidelinks]{hyperref}
\allowdisplaybreaks
\title{Qpu.General.Pi}
\author{Tsvetan Rouschev}
\date{@uuidna/qpu v1.1.1, CC-BY-NC-ND-4.0}
\begin{document}
\maketitle
\noindent\url{https://github.com/uuidna/qpu/blob/v1.1.1/lean/Qpu/General/Pi.lean}

\paragraph{\texttt{pi\_all}}
Every convergent at once: π's floor and its classical rational approximations, the analytic-wing ingredients.
\[ \mathrm{mint}\ \mathrm{coins} - \mathrm{seed} = 3 \land \mathrm{faces} + \mathrm{vertices} = 22 \land \mathrm{rays} = 7 \land \mathrm{plane} \cdot (\mathrm{faces} - \mathrm{seed}) - n \cdot n = 355 \land \mathrm{plane} \cdot \mathrm{hexbit} + \mathrm{seed} = 113 \land (\mathrm{faces} - \mathrm{coins}) \cdot (\mathrm{plane} + \mathrm{coins}) = 360 \land (\mathrm{faces} - \mathrm{coins}) \cdot (\mathrm{faces} + \mathrm{seed}) = 180 \]
\noindent axioms: none; proof: \url{https://github.com/uuidna/qpu/blob/v1.1.1/lean/Qpu/General/Pi.lean\#L23}

\paragraph{\texttt{pi\_archimedes\_den}}
Archimedes' upper bound 22/7: the denominator (7) equals rays.
\[ \mathrm{rays} = 7 \]
\noindent axioms: none; proof: \url{https://github.com/uuidna/qpu/blob/v1.1.1/lean/Qpu/General/Pi.lean\#L12}

\paragraph{\texttt{pi\_archimedes\_num}}
Archimedes' upper bound 22/7: the numerator (22) equals faces + vertices.
\[ \mathrm{faces} + \mathrm{vertices} = 22 \]
\noindent axioms: none; proof: \url{https://github.com/uuidna/qpu/blob/v1.1.1/lean/Qpu/General/Pi.lean\#L10}

\paragraph{\texttt{pi\_half\_turn}}
a half turn in degrees (180), π radians, equals (faces − coins) · (faces + seed).
\[ (\mathrm{faces} - \mathrm{coins}) \cdot (\mathrm{faces} + \mathrm{seed}) = 180 \]
\noindent axioms: none; proof: \url{https://github.com/uuidna/qpu/blob/v1.1.1/lean/Qpu/General/Pi.lean\#L20}

\paragraph{\texttt{pi\_turn\_degrees}}
a full turn in degrees (360) equals (faces − coins) · (plane + coins).
\[ (\mathrm{faces} - \mathrm{coins}) \cdot (\mathrm{plane} + \mathrm{coins}) = 360 \]
\noindent axioms: none; proof: \url{https://github.com/uuidna/qpu/blob/v1.1.1/lean/Qpu/General/Pi.lean\#L18}

\paragraph{\texttt{pi\_zu\_den}}
Zu Chongzhi's milü 355/113: the denominator (113, prime) equals plane · hexbit + seed.
\[ \mathrm{plane} \cdot \mathrm{hexbit} + \mathrm{seed} = 113 \]
\noindent axioms: none; proof: \url{https://github.com/uuidna/qpu/blob/v1.1.1/lean/Qpu/General/Pi.lean\#L16}

\paragraph{\texttt{pi\_zu\_num}}
Zu Chongzhi's milü 355/113: the numerator (355) equals plane · (faces − seed) − n · n.
\[ \mathrm{plane} \cdot (\mathrm{faces} - \mathrm{seed}) - n \cdot n = 355 \]
\noindent axioms: none; proof: \url{https://github.com/uuidna/qpu/blob/v1.1.1/lean/Qpu/General/Pi.lean\#L14}
\end{document}
